\subsection{Governing Equations of Motion}
\label{subsec:governing_eqs}
Transient structural dynamics considers the balance of linear momentum for an elastic body occupying a physical domain $\Omega(\bm{\mu}) \subset \mathbb{R}^d$ ($d=2,3$) dependent on a set of design or operational parameters $\bm{\mu} \in \mathcal{P} \subset \mathbb{R}^{n_p}$. The initial-boundary value problem is governed by:
\begin{equation}
    \rho(\bm{\mu}) \ddot{\bm{u}}(\bm{x}, t; \bm{\mu}) - \nabla \cdot \bm{\sigma}(\bm{u}; \bm{\mu}) = \bm{b}(\bm{x}, t; \bm{\mu}) \quad \text{in } \Omega(\bm{\mu}) \times (0, T],
    \label{eq:momentum_balance}
\end{equation}
subject to Dirichlet and Neumann boundary conditions:
\begin{align}
    \bm{u}(\bm{x}, t; \bm{\mu}) &= \bar{\bm{u}}(\bm{x}, t; \bm{\mu}) \quad \text{on } \Gamma_D(\bm{\mu}) \times (0, T], \\
    \bm{\sigma}(\bm{u}; \bm{\mu}) \cdot \bm{n} &= \bar{\bm{t}}(\bm{x}, t; \bm{\mu}) \quad \text{on } \Gamma_N(\bm{\mu}) \times (0, T],
\end{align}
and initial conditions $\bm{u}(\bm{x}, 0; \bm{\mu}) = \bm{u}_0(\bm{x})$ and $\dot{\bm{u}}(\bm{x}, 0; \bm{\mu}) = \bm{v}_0(\bm{x})$, where $\rho$ is the mass density, $\bm{u}$ is the displacement field, $\bm{\sigma}$ denotes the Cauchy stress tensor, $\bm{b}$ is the body force, and $\bm{n}$ is the outward unit normal vector to the boundary.

\subsection{Linear Elasticity and Weak Form}
\label{subsec:weak_form}
Assuming infinitesimal strain theory, the strain-displacement relation is given by:
\begin{equation}
    \bm{\varepsilon}(\bm{u}) = \frac{1}{2} \left( \nabla \bm{u} + (\nabla \bm{u})^\top \right).
\end{equation}
The constitutive relation follows Hooke's law:
\begin{equation}
    \bm{\sigma}(\bm{u}; \bm{\mu}) = \mathbb{C}(\bm{\mu}) : \bm{\varepsilon}(\bm{u}),
\end{equation}
where $\mathbb{C}(\bm{\mu})$ is the fourth-order elasticity tensor parametrized by Young's modulus $E(\bm{\mu})$ and Poisson's ratio $\nu(\bm{\mu})$.

The variational (weak) formulation consists of finding $\bm{u}(\cdot, t; \bm{\mu}) \in \mathcal{V}$ such that for all test functions $\bm{w} \in \mathcal{V}_0$:
\begin{equation}
    m(\bm{w}, \ddot{\bm{u}}; \bm{\mu}) + c(\bm{w}, \dot{\bm{u}}; \bm{\mu}) + a(\bm{w}, \bm{u}; \bm{\mu}) = f_{\text{ext}}(\bm{w}; \bm{\mu}),
    \label{eq:weak_form}
\end{equation}
where the mass and internal stiffness bilinear forms, alongside the external dynamic load linear functional, are defined as:
\begin{align}
    m(\bm{w}, \ddot{\bm{u}}; \bm{\mu}) &= \int_{\Omega(\bm{\mu})} \rho(\bm{\mu}) \bm{w} \cdot \ddot{\bm{u}} \, \mathrm{d}\Omega, \label{eq:mass_form} \\
    a(\bm{w}, \bm{u}; \bm{\mu}) &= \int_{\Omega(\bm{\mu})} \bm{\varepsilon}(\bm{w}) : \mathbb{C}(\bm{\mu}) : \bm{\varepsilon}(\bm{u}) \, \mathrm{d}\Omega, \label{eq:stiff_form} \\
    f_{\text{ext}}(\bm{w}; \bm{\mu}) &= \int_{\Omega(\bm{\mu})} \bm{w} \cdot \bm{b} \, \mathrm{d}\Omega + \int_{\Gamma_N(\bm{\mu})} \bm{w} \cdot \bar{\bm{t}} \, \mathrm{d}\Gamma, \label{eq:load_form}
\end{align}
and where $\mathcal{V}$ and $\mathcal{V}_0$ denote the trial and test function spaces satisfying inhomogeneous and homogeneous Dirichlet boundary conditions, respectively.

To account for structural dissipation in transient vibration regimes, viscous damping is modeled via the Rayleigh dissipation formulation. In this setting, the dissipation bilinear form is parameterized as a linear combination of inertial and stiffness contributions:
\begin{equation}
    c(\bm{w}, \dot{\bm{u}}; \bm{\mu}) = \alpha_M m(\bm{w}, \dot{\bm{u}}; \bm{\mu}) + \beta_K a(\bm{w}, \dot{\bm{u}}; \bm{\mu}),
    \label{eq:rayleigh_damping}
\end{equation}
where $\alpha_M \ge 0$ and $\beta_K \ge 0$ denote the mass- and stiffness-proportional Rayleigh damping coefficients, calibrated to achieve targeted damping ratios across dominant structural modes.
